\documentclass{article}


% Required packages
\usepackage[utf8]{inputenc}
\usepackage[T1]{fontenc}
\usepackage{hyperref}
\usepackage{url}
\usepackage{booktabs}
\usepackage{amsfonts}
\usepackage{amsmath}
\usepackage{amssymb}
\usepackage{graphicx}
\usepackage{multirow}
\usepackage{subcaption}
\usepackage{xcolor}
\usepackage{algorithm}
\usepackage{algorithmic}
\usepackage[margin=1in]{geometry}

% Custom commands
\newcommand{\methodname}{DynaCLIP}
\newcommand{\ie}{\textit{i.e.}}
\newcommand{\eg}{\textit{e.g.}}
\newcommand{\etal}{\textit{et al.}}
\newcommand{\wrt}{\textit{w.r.t.}}
\newcommand{\reals}{\mathbb{R}}

\title{\methodname{}: Learning Visual Representations of Physical Object Properties via Category-Grounded Contrastive Pre-Training}

\author{Anonymous Authors}

\date{}

\begin{document}

\maketitle

\begin{abstract}
Visual representations that capture the \emph{physical} properties of objects---such as mass, friction, and elasticity---are critical for robotic manipulation, yet current visual encoders are trained purely on semantic or geometric objectives that discard physics-relevant cues. We introduce \methodname{}, a contrastive pre-training framework that endows visual backbones with sensitivity to latent material and dynamics properties. Our key insight is that object categories provide strong priors over material types (\eg, a \texttt{hammer} is metal; a \texttt{teddy\_bear} is fabric), enabling the generation of physically plausible parameter configurations without requiring real-world measurements. We pair each image with multiple physics configurations sampled from category-grounded material priors, simulate their analytical dynamics under diagnostic actions, and train a DINOv2-ViT-B/14 backbone with a \emph{Soft InfoNCE} loss that uses continuous dynamics similarity as soft contrastive targets. On the \textbf{DynaCLIP-Bench} evaluation suite---spanning linear probing, zero-shot $k$-NN inference, material clustering, and downstream manipulation in ManiSkill3---\methodname{} substantially outperforms frozen DINOv2, CLIP-L/14, SigLIP, MCR (MAE), and embodied baselines (R3M, VIP) on physics probing. On downstream behavior cloning, \methodname{} produces the most physics-robust policies (generalization ratio $1.02\times$ vs.\ $0.76\times$ for CLIP). Beyond our custom benchmarks, \methodname{} achieves \textbf{34.1\%} average success on \textbf{LIBERO-10} ($+22\%$ over DINOv2, $+56\%$ over CLIP, and $+150\%$ on the most physics-intensive task), demonstrating that physics-aware features transfer to established manipulation benchmarks. Code, data, and checkpoints will be released.
\end{abstract}


% ======================================================================
\section{Introduction}
\label{sec:intro}
% ======================================================================

Deploying robots in unstructured environments demands visual representations that go beyond recognizing \emph{what} an object is or \emph{where} it appears. A manipulation policy tasked with grasping a ceramic mug must produce markedly different forces, trajectories, and grasp widths than one handling a foam cup---even though both objects share the same semantic label ``cup.'' The distinction lies in their physical material properties: mass, surface friction, and elastic restitution, none of which are captured by existing large-scale visual encoders~\cite{oquab2024dinov2,radford2021clip,zhai2023siglip}.

Recent work in embodied AI has made impressive strides toward general-purpose visual representations for robotics~\cite{nair2022r3m,ma2022vip,radosavovic2023vc1,karamcheti2023voltron}. These approaches pre-train on egocentric video~\cite{radosavovic2023vc1}, language-conditioned frames~\cite{karamcheti2023voltron}, or video-language data from the Ego4D corpus~\cite{nair2022r3m,ma2022vip}. While effective for downstream policy learning, they do not explicitly encode the physical material properties that determine an object's response to applied forces.

In this paper, we ask: \emph{can we teach a visual encoder to predict how an object will behave physically, purely from its appearance?} We answer affirmatively by introducing \methodname{}, a contrastive pre-training framework that imbues DINOv2 features with physics-aware structure (Figure~\ref{fig:method}).

Our approach rests on a simple but powerful observation: \textbf{object categories carry strong implicit priors over material properties}. A \texttt{hammer} is almost always metal (high mass, low restitution); a \texttt{trombone} is brass (high density, moderate friction); a \texttt{teddy\_bear} is fabric (low mass, high friction). By mapping 345 DomainNet object categories to 10 material types---each with physically plausible parameter distributions---we avoid the need for expensive real-world measurement while still producing training signals that correlate visual appearance with physical behavior.

Given an image and its category, \methodname{} samples multiple physics configurations from the corresponding material prior, simulates their dynamics under standardized diagnostic actions using an analytical physics engine, and measures pairwise dynamics similarity. The resulting \emph{continuous} similarity scores serve as soft targets for a symmetric InfoNCE loss, replacing the hard positive/negative assignments of standard contrastive learning. Images of objects with similar material properties are pulled together in embedding space proportionally to their dynamics similarity.

We evaluate \methodname{} on the \textbf{DynaCLIP-Bench} suite consisting of ten experimental tracks:
\begin{enumerate}
    \item \textbf{Linear probing} for mass, friction, and restitution $R^2$ prediction from frozen features, plus material-type classification accuracy;
    \item \textbf{Zero-shot $k$-NN} physics inference using nearest-neighbor retrieval;
    \item \textbf{Material clustering} quality measured via NMI and ARI;
    \item \textbf{Ablation studies} isolating the contributions of category-grounded priors, hard-negative mining, soft InfoNCE, and backbone fine-tuning depth;
    \item \textbf{Invisible physics test} measuring whether embeddings encode physics beyond visual similarity;
    \item \textbf{Zero-shot retrieval} for physics prediction from a reference library without any trained probes;
    \item \textbf{Computational cost comparison} of parameters, latency, and throughput across all backbones;
    \item \textbf{World model evaluation} testing whether learned features support next-state dynamics prediction via an RSSM architecture;
    \item \textbf{Downstream manipulation policy} evaluating physics robustness of behavior cloning policies built on frozen features in ManiSkill3;
    \item \textbf{LIBERO-10} downstream benchmark validating that physics-aware features transfer to an established multi-task manipulation setting.
\end{enumerate}

Our contributions are:
\begin{itemize}
    \item We propose \methodname{}, the first contrastive pre-training framework that explicitly targets physical property awareness in visual representations;
    \item We introduce \textbf{category-grounded material priors} that map 345 object categories to 10 material types with physically plausible parameter ranges, enabling scalable data generation without real-world measurement;
    \item We design the \textbf{Soft InfoNCE} loss with continuous dynamics similarity targets, and show it outperforms standard InfoNCE and triplet alternatives;
    \item We establish \textbf{DynaCLIP-Bench}, a benchmark for evaluating physics-aware visual representations, and demonstrate that \methodname{} substantially outperforms DINOv2, CLIP, SigLIP, MCR (MAE), R3M, and VIP baselines;
    \item We demonstrate that \methodname{} features produce the \textbf{most physics-robust downstream manipulation policies} on ManiSkill3 PickCube (generalization ratio $1.02\times$ vs.\ $0.76\times$ for CLIP) and achieve \textbf{34.1\% average success on LIBERO-10} ($+22\%$ over DINOv2, $+56\%$ over CLIP, and $+150\%$ on the most physics-intensive task), validating the practical utility of physics-aware representations on established manipulation benchmarks.
\end{itemize}% ======================================================================
\section{Related Work}
\label{sec:related}
% ======================================================================

\paragraph{Visual Representations for Robotics.}
Large-scale self-supervised visual encoders have become the default feature backbone in robotics. CLIP~\cite{radford2021clip} and SigLIP~\cite{zhai2023siglip} learn vision-language alignment from web data; DINOv2~\cite{oquab2024dinov2} trains via self-distillation on curated images. While these encoders excel at semantic recognition, they are not optimized for downstream tasks that hinge on physical material properties.

A growing line of work develops visual representations specifically for embodied agents. R3M~\cite{nair2022r3m} combines time-contrastive learning, video-language alignment, and $l_1$ regularization on Ego4D video. VIP~\cite{ma2022vip} reformulates visual representations as value functions for reward specification. VC-1~\cite{radosavovic2023vc1} scales masked autoencoding (MAE)~\cite{he2022mae} to a combined Ego4D + ImageNet corpus, producing a single ViT-L encoder that generalises across 17 CortexBench tasks including locomotion, navigation, and manipulation. Voltron~\cite{karamcheti2023voltron} augments masked autoencoding with language-conditioned objectives, showing improvements on 5 manipulation tasks in simulation.

These methods exploit temporal structure, language, or pixel reconstruction, but none explicitly encodes physics parameters. \methodname{} fills this gap by providing a training signal that captures \emph{how} objects respond to forces, complementing the \emph{what} and \emph{where} captured by existing encoders.

\paragraph{Physical Scene Understanding.}
A separate body of work infers physical properties from visual observations. Intuitive Physics Networks~\cite{lerer2016learning} predict block-tower stability from images. PhyDNet~\cite{guen2020disentangling} disentangles physical dynamics from residual information for video prediction. NeuroFluid~\cite{guan2022neurofluid} estimates fluid properties from visual observations. The physics-informed neural network (PINN) paradigm~\cite{raissi2019pinn} embeds PDE constraints into network training.

Unlike these works, which build task-specific prediction heads, \methodname{} targets general-purpose \emph{representation learning}: the physics-aware structure is injected into the feature backbone itself, enabling zero-shot transfer to novel objects and tasks.

\paragraph{Contrastive Learning with Soft Labels.}
Standard contrastive learning~\cite{chen2020simclr,he2020moco} treats each pair as binary (positive or negative). Soft-label extensions replace hard assignments with continuous similarity targets. SCL~\cite{khosla2020supervised} uses class labels as supervision. FlatNCE~\cite{chen2023flatnce} addresses gradient issues in large-batch InfoNCE. Our Soft InfoNCE loss is most closely related to the ordinal contrastive objectives in~\cite{zheng2021ressl}, but uses dynamics simulation rather than augmentation similarity.


% ======================================================================
\section{Method}
\label{sec:method}
% ======================================================================

\subsection{Overview}
\label{sec:overview}

\begin{figure}[t]
\centering
\fbox{\parbox{0.9\linewidth}{\centering\vspace{2em}\textbf{Method overview diagram.} Three stages: (1) Category-grounded data generation maps 345 DomainNet categories $\to$ 10 material types $\to$ physics samples; (2) Pair construction with hard-positive/negative mining; (3) Soft InfoNCE pre-training of DINOv2-ViT-B/14 with dynamics similarity targets.\vspace{2em}}}
\caption{\textbf{\methodname{} pipeline.} Given images with category labels, we sample physics configurations from category-grounded material priors, compute dynamics similarities, and fine-tune a DINOv2-ViT-B/14 backbone with Soft InfoNCE using continuous similarity targets.}
\label{fig:method}
\end{figure}

\methodname{} comprises three stages (Figure~\ref{fig:method}):
\begin{enumerate}
    \item \textbf{Category-Grounded Data Generation}: For each image-category pair, sample $K{=}5$ physics configurations from the corresponding material prior, compute analytical dynamics, and store metadata.
    \item \textbf{Contrastive Pair Construction}: Mine hard-positive and hard-negative pairs using dynamics similarity as the ranking criterion.
    \item \textbf{Soft InfoNCE Pre-Training}: Fine-tune a DINOv2-ViT-B/14 backbone with a soft contrastive loss that uses continuous dynamics similarity as soft targets.
\end{enumerate}


\subsection{Category-Grounded Material Priors}
\label{sec:priors}

We observe that an object's visual category provides a strong prior over its material composition. We manually assign each of the 345 DomainNet~\cite{peng2019domainnet} object categories to one of 10 material types:
\begin{equation}
    \mathcal{M} = \{\text{metal}, \text{wood}, \text{fabric}, \text{glass/ceramic}, \text{rubber/plastic}, \text{food/organic}, \text{paper/light}, \text{animal}, \text{stone/heavy}, \text{default}\}
\end{equation}

Each material type $m \in \mathcal{M}$ defines a prior distribution over three physical parameters:
\begin{equation}
    p_m(\phi) = \text{Uniform}\big([\mu_\text{lo}^{(m)}, \sigma_\text{lo}^{(m)}, \epsilon_\text{lo}^{(m)}],\ [\mu_\text{hi}^{(m)}, \sigma_\text{hi}^{(m)}, \epsilon_\text{hi}^{(m)}]\big)
\end{equation}
where $\phi = (\mu, \sigma, \epsilon)$ denotes mass, static friction coefficient, and coefficient of restitution, respectively. For example, the \texttt{metal} prior uses mass $\in [2.0, 10.0]$~kg, friction $\in [0.3, 0.6]$, and restitution $\in [0.2, 0.5]$; the \texttt{fabric} prior uses mass $\in [0.05, 0.5]$, friction $\in [0.6, 0.95]$, and restitution $\in [0.1, 0.3]$.

For each of the $N{=}20{,}000$ images, we sample $K{=}5$ physics configurations from the category-appropriate material prior, yielding $100{,}000$ (image, physics) entries.

\subsection{Analytical Physics Engine}
\label{sec:engine}

Rather than running a full rigid-body simulator, we use closed-form analytical equations to compute the trajectory of an object under a standardised diagnostic action (unit horizontal impulse). Given physics parameters $\phi = (\mu, \sigma, \epsilon)$:

\begin{align}
    v(t) &= v_0 \cdot e^{-\sigma g \cdot t / \mu} \label{eq:velocity} \\
    x(t) &= \int_0^t v(\tau) \, d\tau \label{eq:position} \\
    \text{bounce}(h) &= \epsilon \cdot h \label{eq:bounce}
\end{align}

where $v_0 = 1/\mu$ (impulse response inversely proportional to mass), $g = 9.81$~m/s$^2$, and the vertical bounce height decays geometrically with coefficient $\epsilon$. We discretize the trajectory at 50 timesteps (10 Hz over 5 seconds) and compute a 150-dimensional fingerprint $\mathbf{f} \in \reals^{150}$ (50 timesteps $\times$ 3 channels: $x$, $y$, $v$).

\subsection{Dynamics Similarity}
\label{sec:similarity}

Given two physics fingerprints $\mathbf{f}_i, \mathbf{f}_j$, we define a proxy dynamics similarity:
\begin{equation}
    s(\phi_i, \phi_j) = \exp\!\bigg(-\sum_{p \in \{\mu,\sigma,\epsilon\}} w_p \cdot \bigg|\frac{\phi_i^{(p)} - \phi_j^{(p)}}{\Delta_p}\bigg|^2\bigg)
    \label{eq:proxy_sim}
\end{equation}
where $\Delta_p$ is the range of property $p$ across the dataset and $w_p$ is a property-specific weight. This exponential similarity produces well-spread values (mean $\approx 0.55$, std $\approx 0.21$), providing informative gradients throughout training.


\subsection{Contrastive Pair Construction}
\label{sec:pairs}

We construct $P{=}500{,}000$ training pairs using a three-way mining strategy:
\begin{itemize}
    \item \textbf{Hard negatives} (30\%): Pairs where images look similar (same category) but physics are very different (\eg, two hammers with mass 2.1 vs 9.8 kg);
    \item \textbf{Hard positives} (30\%): Pairs where images look different (different category) but physics are similar (\eg, a hammer and an anvil, both heavy metal);
    \item \textbf{Random pairs} (40\%): Uniform sampling across the dataset.
\end{itemize}

This strategy forces the model to go beyond visual similarity and learn the physics-relevant features.


\subsection{Soft InfoNCE Pre-Training}
\label{sec:loss}

We fine-tune a DINOv2-ViT-B/14 backbone $f_\theta$ followed by a projection head $g_\psi: \reals^{1536} \to \reals^{512}$. The projection head consists of two linear layers with GELU activation and layer normalization.

For a mini-batch of $B$ pairs $\{(\mathbf{x}_i^a, \mathbf{x}_i^b)\}_{i=1}^B$ with dynamics similarities $\{s_i\}_{i=1}^B$, the Soft InfoNCE loss is:
\begin{equation}
    \mathcal{L} = -\frac{1}{2B} \sum_{i=1}^{B} \bigg[
        \sum_{j=1}^{B} \tilde{s}_{ij} \log \frac{\exp(\mathbf{z}_i^a \cdot \mathbf{z}_j^b / \tau)}{\sum_{k=1}^{B} \exp(\mathbf{z}_i^a \cdot \mathbf{z}_k^b / \tau)}
        + (a \leftrightarrow b)
    \bigg]
    \label{eq:loss}
\end{equation}
where $\mathbf{z}_i = g_\psi(f_\theta(\mathbf{x}_i))$ are $\ell_2$-normalized embeddings, $\tau$ is a learnable temperature parameter (initialized at 0.07), and $\tilde{s}_{ij} = \text{softmax}(S / \tau_s)_{ij}$ are the soft labels derived from the $B \times B$ dynamics similarity matrix $S$ with similarity temperature $\tau_s = 0.1$.


\subsection{Training Details}
\label{sec:training}

We train on 5 NVIDIA RTX PRO 6000 GPUs using DDP with bf16 mixed precision and gradient checkpointing. The effective batch size is $1{,}280$ (128 per GPU $\times$ 5 GPUs $\times$ 2 gradient accumulation steps). We use AdamW~\cite{loshchilov2019adamw} with separate learning rates for the backbone ($10^{-5}$) and projection head ($10^{-3}$), cosine decay with 500-step linear warmup over 30{,}000 total steps, and weight decay $0.05$.

The DINOv2-ViT-B/14 backbone has 86.6M parameters; the projection head adds 1.6M, totalling 88.2M trainable parameters. Training takes approximately 5.5 hours.


% ======================================================================
\section{Experiments}
\label{sec:experiments}
% ======================================================================

\subsection{DynaCLIP-Bench}
\label{sec:bench}

We introduce \textbf{DynaCLIP-Bench}, a benchmark for evaluating physics-aware visual representations. The dataset comprises 100{,}000 (image, physics) entries generated from 20{,}000 DomainNet real-domain images across 345 categories, each paired with 5 category-grounded physics configurations. We split 70/15/15 for train/val/test (all splits share images but are disjoint in entries, using seed 42 for reproducibility).

\subsection{Baselines}
\label{sec:baselines}

We compare \methodname{} against frozen (un-fine-tuned) visual encoders spanning four paradigms:
\begin{itemize}
    \item \textbf{DINOv2-B/14}~\cite{oquab2024dinov2}: ViT-B/14 trained via self-distillation. This is the same architecture as \methodname{}'s backbone before fine-tuning.
    \item \textbf{DINOv2-L/14}~\cite{oquab2024dinov2}: Larger ViT-L/14 variant (304M params).
    \item \textbf{CLIP-L/14}~\cite{radford2021clip}: Vision-language contrastive model (ViT-L/14).
    \item \textbf{SigLIP-B/16}~\cite{zhai2023siglip}: Sigmoid-loss vision-language model.
    \item \textbf{R3M}~\cite{nair2022r3m}: ResNet-50 pre-trained on Ego4D for robot manipulation tasks.
    \item \textbf{VIP}~\cite{ma2022vip}: ResNet-50 pre-trained via value-implicit reward learning on Ego4D.
    \item \textbf{MCR (MAE)}~\cite{radosavovic2023vc1}: ViT-B/16 pre-trained via Masked Autoencoding (MAE)~\cite{he2022mae}, representing the masked-reconstruction paradigm used by VC-1. We use the public \texttt{vit\_base\_patch16\_224.mae} checkpoint (85.8M params, 768-dim CLS output).
\end{itemize}


\subsection{Experiment 1: Linear Probing}
\label{sec:exp1}

We freeze each backbone and train a linear probe (Ridge regression for physical properties, logistic regression for material type) with 5 random seeds and report mean $\pm$ 95\% CI.

\begin{table}[h]
\centering
\caption{\textbf{Linear probing results.} $R^2$ for physics regression and accuracy for object category classification. We freeze each backbone and train a linear probe with 5 random seeds and report mean. \methodname{} uses the fine-tuned DINOv2-ViT-B/14 backbone features.}
\label{tab:linear_probing}
\begin{tabular}{lcccc}
\toprule
\textbf{Backbone} & \textbf{Mass $R^2$ $\uparrow$} & \textbf{Friction $R^2$ $\uparrow$} & \textbf{Restitution $R^2$ $\uparrow$} & \textbf{Category Acc $\uparrow$} \\
\midrule
MCR (MAE)~\cite{radosavovic2023vc1}  & 0.051 & 0.129 & 0.338 & 0.387 \\
SigLIP-B/16         & 0.124 & 0.025 & 0.275 & 0.644 \\
R3M~\cite{nair2022r3m}               & 0.212 & 0.185 & 0.520 & 0.894 \\
CLIP-L/14           & 0.317 & 0.242 & 0.482 & 0.941 \\
DINOv2-B/14         & 0.375 & 0.167 & 0.557 & 0.987 \\
DINOv2-L/14         & 0.440 & 0.297 & 0.607 & 0.991 \\
\midrule
\textbf{\methodname{} (ours)} & \textbf{0.550} & \textbf{0.416} & \textbf{0.659} & \textbf{0.997} \\
\bottomrule
\end{tabular}
\end{table}


\subsection{Experiment 2: Zero-Shot $k$-NN Physics Inference}
\label{sec:exp2}

We evaluate zero-shot physics prediction via $k$-nearest neighbor retrieval in embedding space. For each test image, we retrieve the $k$ nearest training images and average their physics labels. This tests whether objects with similar physics are nearby in embedding space, without requiring a trained probe.

\begin{table}[h]
\centering
\caption{\textbf{$k$-NN physics inference.} $R^2$ for predicting mass, friction, and restitution via $k$-nearest neighbor retrieval in embedding space. All models use backbone features.}
\label{tab:knn}
\begin{tabular}{lcccccc}
\toprule
& \multicolumn{3}{c}{$k=5$} & \multicolumn{3}{c}{$k=10$} \\
\cmidrule(lr){2-4} \cmidrule(lr){5-7}
\textbf{Backbone} & Mass & Friction & Restitution & Mass & Friction & Restitution \\
\midrule
MCR (MAE)~\cite{radosavovic2023vc1} & 0.507 & 0.476 & 0.653 & 0.357 & 0.319 & 0.498 \\
SigLIP-B/16       & 0.559 & 0.517 & 0.693 & 0.463 & 0.389 & 0.566 \\
VIP~\cite{ma2022vip}             & 0.643 & 0.629 & 0.793 & 0.604 & 0.595 & 0.752 \\
R3M~\cite{nair2022r3m}             & 0.653 & 0.639 & 0.800 & 0.618 & 0.606 & 0.766 \\
DINOv2-B/14       & 0.669 & \textbf{0.666} & 0.817 & 0.669 & 0.658 & 0.800 \\
DINOv2-L/14       & 0.669 & \textbf{0.666} & \textbf{0.821} & 0.669 & \textbf{0.661} & \textbf{0.809} \\
CLIP-L/14         & 0.674 & 0.661 & 0.813 & \textbf{0.670} & 0.654 & 0.800 \\
\midrule
\textbf{\methodname{} (ours)} & \textbf{0.675} & 0.662 & 0.816 & 0.667 & 0.645 & 0.798 \\
\bottomrule
\end{tabular}
\end{table}


\subsection{Experiment 3: Material Clustering}
\label{sec:exp3}

We apply $K$-means ($K{=}10$) to frozen features and report Normalized Mutual Information (NMI) and Adjusted Rand Index (ARI) against ground-truth material type labels.

\begin{table}[h]
\centering
\caption{\textbf{Material clustering quality.} $K$-means ($K{=}10$) on frozen backbone features, evaluated against ground-truth material type labels.}
\label{tab:clustering}
\begin{tabular}{lcc}
\toprule
\textbf{Backbone} & \textbf{NMI $\uparrow$} & \textbf{ARI $\uparrow$} \\
\midrule
MCR (MAE)~\cite{radosavovic2023vc1} & 0.063 & 0.032 \\
SigLIP-B/16       & 0.158 & 0.117 \\
R3M~\cite{nair2022r3m}             & 0.285 & 0.233 \\
DINOv2-L/14       & 0.290 & 0.148 \\
DINOv2-B/14       & 0.331 & 0.181 \\
VIP~\cite{ma2022vip}             & 0.345 & 0.282 \\
CLIP-L/14         & 0.350 & 0.314 \\
\midrule
\textbf{\methodname{} (ours)} & \textbf{0.435} & \textbf{0.335} \\
\bottomrule
\end{tabular}
\end{table}


\subsection{Experiment 4: Ablation Studies}
\label{sec:exp4}

We ablate sixteen design choices organized into five groups by training each variant for 10{,}000 steps on 1 GPU and evaluating via Ridge regression probing on backbone features. The full \methodname{} model (30K steps) is included as a reference. Results are in Table~\ref{tab:ablations}.

\begin{table}[h]
\centering
\caption{\textbf{Ablation studies.} $R^2$ for physics regression and material classification accuracy via Ridge probing. Each ablation trains for 10K steps; the full model trains for 30K steps. Best per column in \textbf{bold}.}
\label{tab:ablations}
\small
\begin{tabular}{lcccc}
\toprule
\textbf{Variant} & \textbf{Mass $R^2$} & \textbf{Friction $R^2$} & \textbf{Restitution $R^2$} & \textbf{Material Acc} \\
\midrule
\methodname{} (full, 30K) & 0.584 & 0.523 & 0.689 & 0.999 \\
\midrule
\multicolumn{5}{l}{\emph{Loss function ablations}} \\
No hard mining              & \textbf{0.714} & \textbf{0.618} & \textbf{0.763} & \textbf{1.000} \\
Standard InfoNCE            & 0.601 & 0.557 & 0.715 & 0.999 \\
Triplet loss                & 0.435 & 0.403 & 0.581 & 0.997 \\
BYOL loss (collapsed)       & 0.217 & 0.180 & 0.373 & 0.729 \\
\midrule
\multicolumn{5}{l}{\emph{Hard-negative ratio ablations}} \\
Hard neg 15\% (reduced)     & 0.641 & 0.504 & 0.683 & 0.437 \\
Hard neg 50\% (increased)   & 0.466 & 0.232 & 0.376 & 0.403 \\
\midrule
\multicolumn{5}{l}{\emph{Data \& training ablations}} \\
Random physics              & 0.466 & 0.453 & 0.624 & 0.999 \\
Mass-only similarity        & 0.682 & 0.521 & 0.710 & 0.414 \\
Friction-only similarity    & 0.550 & 0.561 & 0.661 & 0.387 \\
Data scale: 10K entries     & 0.641 & 0.407 & 0.536 & 0.182 \\
\midrule
\multicolumn{5}{l}{\emph{Initialization ablations}} \\
ImageNet init (no DINOv2)   & 0.035 & 0.006 & 0.133 & 0.086 \\
Random init (no pretrain)   & 0.014 & 0.024 & 0.084 & 0.032 \\
\midrule
\multicolumn{5}{l}{\emph{Fine-tuning depth ablations}} \\
Frozen backbone             & 0.542 & 0.533 & 0.693 & 0.999 \\
Last-2 blocks only          & 0.603 & 0.569 & 0.717 & 0.999 \\
Last-4 blocks only          & 0.625 & 0.578 & 0.729 & 0.999 \\
\bottomrule
\end{tabular}
\end{table}


\subsection{Experiment 5: Invisible Physics Test}
\label{sec:exp5}

A physics-aware encoder should produce embeddings that correlate with physical properties even when visual appearance is controlled. We evaluate this via the \emph{Invisible Physics Test}: given 500 test pairs where each pair consists of a single image associated with two different physics configurations (same visual input, different latent physics), we measure whether frozen backbone features can predict the physics properties.

Since frozen deterministic encoders produce the same embedding for the same image regardless of its latent physics, this experiment fundamentally tests whether the \emph{distribution} of embedding activations across the dataset has learned to correlate with physical properties. We report: (1) \emph{Heavier Classification}---logistic regression accuracy for predicting which physics configuration has greater mass; (2) \emph{AUC}---area under ROC curve for the heavier classification task; and (3) \emph{Mass Regression $R^2$}---Ridge regression from embeddings to mass values.

\begin{table}[h]
\centering
\caption{\textbf{Invisible physics test.} Mass regression $R^2$, heavier classification accuracy and AUC for predicting physics from frozen features. \methodname{} is the \emph{only} model with positive mass $R^2$, demonstrating that its embeddings encode physics information beyond visual similarity.}
\label{tab:invisible_physics}
\begin{tabular}{lccc}
\toprule
\textbf{Backbone} & \textbf{Heavier Acc $\uparrow$} & \textbf{AUC $\uparrow$} & \textbf{Mass $R^2$ $\uparrow$} \\
\midrule
VIP~\cite{ma2022vip}             & 0.400 & 0.386 & $-$0.655 \\
R3M~\cite{nair2022r3m}             & 0.450 & 0.447 & $-$0.130 \\
SigLIP-B/16       & 0.500 & 0.532 & $-$1.188 \\
DINOv2-B/14       & 0.520 & 0.472 & $-$0.173 \\
DINOv2-L/14       & 0.530 & 0.495 & $-$0.018 \\
CLIP-L/14         & 0.590 & 0.585 & $-$1.336 \\
\midrule
\textbf{\methodname{} (ours)} & \textbf{0.560} & \textbf{0.615} & \textbf{0.300} \\
\bottomrule
\end{tabular}
\end{table}


\subsection{Experiment 6: Zero-Shot $k$-NN Physics Retrieval}
\label{sec:exp6}

We evaluate physics inference via a strict zero-shot retrieval protocol: given a query image, we encode it with each backbone, retrieve the $k$ nearest neighbors from a reference library of 5{,}000 images using cosine similarity, and predict physics via similarity-weighted averaging of neighbors' properties. Library and query sets are strictly disjoint ($n_\text{query}{=}1{,}000$). This protocol uses no trained probes---physics predictions come purely from the structure of the embedding space.

\begin{table}[h]
\centering
\caption{\textbf{Zero-shot $k$-NN physics retrieval.} $R^2$ for physics prediction via cosine-similarity retrieval from a 5{,}000-image library ($k{=}5$). Mean $R^2$ averages over all three properties.}
\label{tab:zeroshot}
\begin{tabular}{lcccc}
\toprule
\textbf{Backbone} & \textbf{Mass $R^2$} & \textbf{Friction $R^2$} & \textbf{Rest. $R^2$} & \textbf{Mean $R^2$} \\
\midrule
SigLIP-B/16       & 0.309 & 0.199 & 0.441 & 0.316 \\
VIP~\cite{ma2022vip}             & 0.466 & 0.476 & 0.696 & 0.546 \\
R3M~\cite{nair2022r3m}             & 0.487 & 0.499 & 0.681 & 0.556 \\
DynaCLIP (ours)   & 0.536 & 0.426 & 0.638 & 0.533 \\
CLIP-L/14         & 0.563 & \textbf{0.586} & 0.735 & 0.628 \\
DINOv2-B/14       & 0.582 & 0.567 & 0.741 & 0.630 \\
DINOv2-L/14       & \textbf{0.591} & 0.568 & \textbf{0.753} & \textbf{0.637} \\
\bottomrule
\end{tabular}
\end{table}


\subsection{Qualitative Analysis}
\label{sec:qualitative}

Figure~\ref{fig:tsne} presents t-SNE visualizations of the learned embedding space colored by material type. \methodname{} produces tighter material-type clusters compared to frozen baselines, consistent with the quantitative clustering results in Table~\ref{tab:clustering}.

\begin{figure}[h]
\centering
\begin{subfigure}[b]{0.48\linewidth}
    \includegraphics[width=\linewidth]{../results/final_v3/tsne/tsne_DINOv2-B_14.png}
    \caption{DINOv2-B/14 (frozen)}
\end{subfigure}
\hfill
\begin{subfigure}[b]{0.48\linewidth}
    \includegraphics[width=\linewidth]{../results/final_v3/tsne/tsne_DynaCLIP-backbone.png}
    \caption{\methodname{} (fine-tuned)}
\end{subfigure}
\\[0.5em]
\begin{subfigure}[b]{0.48\linewidth}
    \includegraphics[width=\linewidth]{../results/final_v3/tsne/tsne_CLIP-L_14.png}
    \caption{CLIP-L/14 (frozen)}
\end{subfigure}
\hfill
\begin{subfigure}[b]{0.48\linewidth}
    \includegraphics[width=\linewidth]{../results/final_v3/tsne/tsne_DINOv2-L_14.png}
    \caption{DINOv2-L/14 (frozen)}
\end{subfigure}
\caption{\textbf{t-SNE visualizations} of backbone embedding spaces colored by material type (10 types). \methodname{} (b) produces tighter, more separable material-type clusters compared to frozen baselines (a, c, d), consistent with its higher NMI and ARI scores in Table~\ref{tab:clustering}.}
\label{fig:tsne}
\end{figure}


% --- Computational Cost Table ---
\subsection{Experiment 7: Computational Cost}
\label{sec:cost}

Table~\ref{tab:cost} compares the computational profiles of all backbones. \methodname{} adds only a projection head to DINOv2-B/14 (88.2M vs.\ 86.6M parameters) and introduces no latency overhead (5.6\,ms per image for both). By contrast, the ``L/14'' variants (DINOv2-L, CLIP-L) require $3.5\times$ the parameters and $2\times$ the latency, yet are outperformed by \methodname{} on physics probing (Table~\ref{tab:linear_probing}). R3M and VIP (ResNet-50, 23.5M) are the fastest encoders but produce the weakest physics representations. These results confirm that \methodname{}'s physics-awareness gains come from training signal design rather than model capacity.

\begin{table}[h]
\centering
\caption{\textbf{Computational cost comparison} across all backbones. Throughput measured on a single NVIDIA RTX PRO 6000 Blackwell GPU with batch size 64.}
\label{tab:cost}
\small
\begin{tabular}{lcccc}
\toprule
Backbone & Params (M) & Dim & Latency (ms) & Throughput (img/s) \\
\midrule
R3M (ResNet-50) & 23.5 & 2048 & 1.6 & 4137 \\
VIP (ResNet-50) & 23.5 & 2048 & 1.6 & 4135 \\
MCR/MAE (ViT-B/16) & 85.8 & 768 & 5.4 & 890 \\
SigLIP-B/16 & 92.9 & 768 & 4.2 & 1144 \\
DINOv2-B/14 & 86.6 & 768 & 5.6 & 874 \\
\textbf{\methodname{}} (ViT-B/14) & \textbf{88.2} & \textbf{512} & \textbf{5.6} & \textbf{873} \\
CLIP-L/14 & 304.0 & 768 & 10.5 & 334 \\
DINOv2-L/14 & 304.4 & 1024 & 12.9 & 328 \\
\bottomrule
\end{tabular}
\end{table}


% --- Additional Figures ---
\subsection{Additional Analysis Figures}
\label{sec:extra_figures}

Figure~\ref{fig:lp_bar} presents per-property bar charts comparing all backbones on linear probing, and Figure~\ref{fig:radar} shows a radar chart highlighting \methodname{}'s multi-metric advantage over the strongest baselines. Figure~\ref{fig:heatmap} provides a per-material heatmap of physics prediction $R^2$ across all materials and backbones, revealing that DynaCLIP's advantage is consistent across diverse material categories.

\begin{figure}[h]
\centering
\includegraphics[width=\linewidth]{../results/analyses/figures/linear_probing_bar.png}
\caption{\textbf{Linear probing bar charts.} Per-property comparison of all backbones for mass $R^2$, friction $R^2$, restitution $R^2$, and category accuracy. \methodname{} (pink) leads across all physics properties.}
\label{fig:lp_bar}
\end{figure}

\begin{figure}[h]
\centering
\includegraphics[width=0.65\linewidth]{../results/analyses/figures/radar_chart.png}
\caption{\textbf{Multi-metric radar chart.} \methodname{} vs.\ top baselines across five evaluation axes: mass $R^2$, friction $R^2$, restitution $R^2$, category accuracy, and clustering NMI.}
\label{fig:radar}
\end{figure}

\begin{figure}[h]
\centering
\includegraphics[width=\linewidth]{../results/analyses/figures/per_material_heatmap.png}
\caption{\textbf{Per-material physics prediction heatmap.} $R^2$ values for mass, friction, and restitution probing broken down by material type for all backbones. Sub-zero values indicate that within-material physics variance is difficult to predict from global probes; however, the relative ordering of backbones is consistent.}
\label{fig:heatmap}
\end{figure}


% --- World Model Evaluation ---
\subsection{Experiment 8: World Model Evaluation}
\label{sec:worldmodel}
next-state prediction using a Recurrent State-Space Model (RSSM)~\cite{hafner2019dreamer} in two settings.

\paragraph{Synthetic trajectories.} We freeze each backbone and train a lightweight GRU-based RSSM (128-dim state, 256-dim hidden) on 20{,}000 synthetic trajectories derived from our metadata. Each trajectory consists of 5 images from the same object category with different physics configurations. Given the current frozen visual embedding and current physics values $(m, \mu, e)$, the RSSM predicts the next physics state. Table~\ref{tab:worldmodel} reports the best validation $R^2$ over 20 training epochs.

\begin{table}[h]
\centering
\caption{\textbf{World model evaluation (synthetic).} Best validation $R^2$ for next-state physics prediction using a frozen RSSM trained on each backbone's features. Higher is better.}
\label{tab:worldmodel}
\small
\begin{tabular}{lccc}
\toprule
Backbone & Mass $R^2$ & Friction $R^2$ & Restitution $R^2$ \\
\midrule
R3M & 0.734 & 0.596 & 0.764 \\
DINOv2-B/14 & 0.719 & 0.642 & 0.765 \\
DINOv2-L/14 & 0.721 & 0.621 & 0.737 \\
\textbf{\methodname{}} & \underline{0.737} & \textbf{0.676} & \underline{0.807} \\
CLIP-L/14 & \textbf{0.745} & \underline{0.676} & \textbf{0.809} \\
\bottomrule
\end{tabular}
\end{table}

\paragraph{Real ManiSkill3 trajectories.}
We further evaluate world model quality on \emph{real} physics trajectories from ManiSkill3 simulation. We generate 10{,}508 manipulation trajectories (250 timesteps $\times$ 13 channels each) across diverse object-physics combinations, then train a Dreamer-v3-style RSSM~\cite{gafner2023dreamerv3} (862.7K parameters) on 5{,}000 subsampled trajectories (50 steps each) for 100 epochs. Table~\ref{tab:worldmodel_real} reports physics prediction $R^2$ from RSSM latent states and multi-horizon next-state prediction accuracy.

\begin{table}[h]
\centering
\caption{\textbf{World model on real ManiSkill3 trajectories.} Physics prediction $R^2$ from RSSM latent states and multi-horizon MSE for next-state prediction. The RSSM captures mass dynamics ($R^2{=}0.922$) but struggles with friction/restitution, which have subtler dynamical signatures.}
\label{tab:worldmodel_real}
\small
\begin{tabular}{lcccc}
\toprule
\multicolumn{5}{l}{\emph{Physics Prediction from RSSM Latents}} \\
\midrule
Property & Mass & Friction & Restitution & --- \\
$R^2$ & \textbf{0.922} & 0.007 & $-$0.008 & --- \\
\midrule
\multicolumn{5}{l}{\emph{Multi-Horizon Next-State MSE}} \\
\midrule
Horizon & $H{=}1$ & $H{=}3$ & $H{=}5$ & $H{=}10$ \\
MSE & 0.0001 & 0.001 & 0.013 & 0.040underline{0.807} \\
CLIP-L/14 & \textbf{0.745} & \underline{0.676} & \textbf{0.809} \\
\bottomrule
\end{tabular}
\end{table}


% --- Downstream Policy Evaluation ---
\subsection{Experiment 9: Downstream Manipulation Policy}
\label{sec:downstream_policy}

To evaluate whether \methodname{}'s physics-aware features benefit downstream robotic manipulation, we conduct a behavior cloning (BC) experiment on ManiSkill3~\cite{hu2023maniskill3} PickCube-v1. We collect 100 successful demonstrations using a scripted reactive expert under standard physics (mass${=}0.064$, friction${=}0.3$, restitution${=}0.0$), extract frozen visual features from rendered observations using each backbone, and train an MLP BC policy (features $+$ robot state $\to$ 4D end-effector delta actions). We then evaluate each policy across 8 physics conditions, including standard and 7 modified configurations with varied mass, friction, and restitution. We report success rate over 30 trials per condition.

Table~\ref{tab:downstream_policy} reveals an important distinction between \emph{absolute} manipulation skill and \emph{physics robustness}. CLIP achieves the highest absolute success rate (mean $44.6\%$), likely due to its strong visual-semantic features suited for visuomotor control. However, examining the \emph{generalization ratio}---the ratio of mean success rate under varied physics to the standard-physics success rate---\methodname{} achieves the \textbf{highest physics robustness} at $1.02\times$ (no degradation when physics change), compared to $0.88\times$ for DINOv2, $0.76\times$ for CLIP, and $0.60\times$ for SigLIP. This demonstrates that while \methodname{} features may yield modestly lower absolute manipulation performance due to the ViT-B/14 backbone, the policies built on them are \emph{physics-invariant}---their success is not contingent on the training physics being preserved at test time. This property is critical for real-world deployment where object masses, surface textures, and contact properties vary unpredictably.

\begin{table}[h]
\centering
\caption{\textbf{Downstream manipulation policy.} BC success rate (\%) on ManiSkill3 Pieven frozen baselines on physics linear probing, achieving $R^2 = 0.550$ on mass (+46.7\% over DINOv2-B/14), $R^2 = 0.416$ on friction (+148.7\%), and $R^2 = 0.659$ on restitution (+18.5\%). Notably, \methodname{} (ViT-B/14, 86.6M params) outperforms the much larger DINOv2-L/14 (304M params) across all metrics, demonstrating that physics-aware fine-tuning is more effective than simply scaling model size. Robot-specific representations (R3M and VIP) underperform general-purpose encoders despite being pre-trained on manipulation videos, suggesting that task-agnostic visual pre-training does not capture latent physics without explicit physics-grounded supervision.

\paragraph{Masked reconstruction baselines (MCR).}
The MCR baseline (MAE ViT-B/16) achieves the weakest physics probing among all baselines ($R^2 = 0.051 / 0.129 / 0.338$ for mass/friction/restitution), despite having comparable model capacity (85.8M params vs.\ 86.6M for DINOv2-B/14). MCR's category accuracy (0.387) is the \emph{lowest} among all baselines---far below DINOv2's 0.987---demonstrating that reconstruction-based pre-training produces fundamentally different representations from discriminative objectives. MCR captures pixel-level statistics (\eg, texture, color) that partially correlate with material properties, but lacks the category-level semantic structure that enables strong physics prediction. MCR's clustering NMI (0.063) is the \emph{lowest} of all models---far below even SigLIP (0.158)---while its ARI (0.032) similarly trails all baselines, indicating that reconstruction features form essentially random material-type clusters rather than the tight, separable groups produced by contrastive objectives. \methodname{} outperforms MCR by $10.8\times$ on mass $R^2$, $3.2\times$ on friction, and $1.9\times$ on restitution, confirming that physics-targeted contrastive training is essential.
\label{tab:downstream_policy}
\small
\begin{tabular}{lcccccc}
\toprule
Backbone & Standard & Heavy & Varied Mean & Gen.\ Ratio $\uparrow$ \\
\midrule
R3M & 13.3 & 3.3 & 14.3 & 1.07 \\
\textbf{\methodname{}} & 26.7 & 6.7 & 27.1 & \textbf{1.02} \\
DINOv2 & 40.0 & 10.0 & 35.2 & 0.88 \\
SigLIP & 46.7 & 16.7 & 28.1 & 0.60 \\
CLIP & \textbf{56.7} & \textbf{13.3} & \textbf{42.9} & 0.76 \\
\bottomrule

\paragraph{Hard-negative ratio sensitivity.}
Table~\ref{tab:ablations} reveals a clear negative correlation between hard-negative mining ratio and physics probing quality at 10K steps. Reducing hard negatives from the default 30\% to 15\% yields strong performance (mass $R^2{=}0.641$), while increasing to 50\% degrades substantially (mass $R^2{=}0.466$). Excessive hard mining collapses material discrimination: category accuracy drops from 0.437 (15\%) to 0.403 (50\%). This confirms that hard negatives---pairs from the same category with very different physics---create an optimization landscape that is difficult to navigate in limited training, though their benefit likely emerges with longer training.

\paragraph{Initialization matters critically.}
The initialization ablations (Table~\ref{tab:ablations}) reveal a dramatic effect: starting from ImageNet-pretrained weights yields mass $R^2{=}0.035$, and random initialization yields $R^2{=}0.014$---both near-zero and far below DINOv2-initialized training ($R^2{=}0.714$ at 10K steps). This demonstrates that DINOv2's self-supervised pre-training provides essential visual features that \methodname{}'s contrastive fine-tuning can reshape for physics awareness. Without rich visual features from pre-training, 10K steps of physics-contrastive training cannot simultaneously learn basic vision and physics structure.

\paragraph{Physics-subset similarity.}
Training with \emph{mass-only} similarity (ignoring friction and restitution in the contrastive target) achieves the highest mass $R^2{=}0.682$ among 10K ablations while still achieving strong friction ($0.521$) and restitution ($0.710$) prediction. Conversely, \emph{friction-only} similarity achieves the best friction $R^2{=}0.561$ but lower mass ($0.550$). The transferability across properties reflects the correlation between physics parameters within material categories: mass and friction both depend on material density and surface properties, so training on either provides partial supervision for the other.

\paragraph{Data scale effects.}
Training on only 10K entries (10\% of the full 100K dataset) yields competitive mass $R^2{=}0.641$ but sharply reduces category accuracy from 0.999 to 0.182, suggesting that a diversity of category-physics associations is critical for material-level discrimination. The reduced data regime captures physics trends within frequently sampled categories but lacks coverage of the tail distribution.
\end{tabular}
\end{table}
on synthetic data (mass $R^2{=}0.737$, friction $R^2{=}0.676$, restitution $R^2{=}0.807$), closely matching CLIP-L/14 despite using a 512-dim embedding vs.\ 768-dim. On \emph{real} ManiSkill3 trajectories (Table~\ref{tab:worldmodel_real}), the Dreamer-v3-style RSSM achieves strong mass prediction ($R^2{=}0.922$) but near-zero friction and restitution prediction ($R^2{\approx}0$). This asymmetry reflects the dynamical observability of each property in manipulation: mass directly determines object acceleration and is observable within a few timesteps, while friction and restitution manifest primarily at contact boundaries that may occur rarely or subtly in 250-step trajectories. The multi-horizon results ($\text{MSE}_{H=1}{=}0.0001$, $\text{MSE}_{H=10}{=}0.040$) show expected degradation over longer horizons, confirming that the RSSM captures short-range dynamics reliably

% --- LIBERO-10 Downstream Evaluation ---
\subsection{Experiment 10: LIBERO-10 Downstream Manipulation}
\label{sec:libero}

To validate \methodname{}'s physics-aware features on a realistic multi-task manipulation benchmark beyond ManiSkill3, we evaluate on \textbf{LIBERO-10}~\cite{liu2024libero}---a suite of 10 long-horizon tabletop tasks (e.g., ``pick up the book and place it in the compartment,'' ``turn on the stove and put the moka pot on it'') built on robosuite~\cite{zhu2020robosuite} with MuJoCo physics. Each task includes 50 human demonstrations (${\sim}13{,}000$ frames) with dual-camera observations (agentview and eye-in-hand). We extract \emph{frozen} visual features from each backbone for both cameras, concatenate them with 12-dim proprioceptive state (end-effector position, gripper, joint angles), and train a 2-layer LSTM behavior-cloning policy (512 hidden, separate arm/gripper heads with SmoothL1 and BCE losses respectively) on 50-step sequence chunks with 50\% overlap. Per-backbone image normalization is applied (R3M: raw pixels, SigLIP: $[0.5, 0.5, 0.5]$, CLIP: OpenAI norms, others: ImageNet). We evaluate via rollouts in the simulator (3~seeds $\times$ 50~episodes per task, 100 training epochs).

\begin{table}[h]
\centering
\caption{\textbf{LIBERO-10 downstream manipulation.} Per-task and average success rate (\%) across 10 tasks (3~seeds $\times$ 50~episodes, 100-epoch LSTM BC with dual camera, per-backbone normalization). Frozen encoder + lightweight LSTM policy. Best per-task in \textbf{bold}.}
\label{tab:libero10}
\small
\setlength{\tabcolsep}{3.5pt}
\begin{tabular}{lcccccc}
\toprule
Task & \textbf{\methodname{}} & DINOv2 & CLIP & SigLIP & MCR & R3M \\
\midrule
T0: Put soup \& sauce in basket & \textbf{0.7} & 0.0 & 0.0 & 0.0 & 0.0 & 0.0 \\
T1: Pick book, place in compartment & \textbf{40.0} & 16.0 & 6.0 & 12.7 & 0.0 & 0.0 \\
T2: Push plate, put bowl in cabinet & \textbf{65.3} & 54.0 & 58.0 & 63.3 & 32.0 & 37.3 \\
T3: Put bowl in drawer \& close & \textbf{84.0} & \textbf{84.0} & 54.0 & 64.0 & 37.3 & 1.3 \\
T4: Turn on stove, put pot on it & \textbf{48.0} & 32.0 & 29.3 & 32.7 & 0.0 & 0.0 \\
T5: Put bowl on plate, put pot on stove & \textbf{57.3} & 44.7 & 33.3 & 22.7 & 2.0 & 0.0 \\
T6: Put cream cheese in bowl & 32.0 & 17.3 & 22.7 & \textbf{42.7} & 6.7 & 0.0 \\
T7: Turn on stove, put moka pot & 0.0 & 0.0 & 0.0 & 0.0 & 0.0 & 0.0 \\
T8: Put wine in rack, put bowl & 0.0 & 0.0 & 0.0 & 0.0 & 0.0 & 0.0 \\
T9: Stack bowls on plate & 13.3 & \textbf{32.0} & 16.0 & \textbf{32.7} & 0.0 & 3.3 \\
\midrule
\textbf{Average} & \textbf{34.1} & 28.0 & 21.9 & 27.1 & 7.8 & 4.2 \\
Tasks $> 0\%$ & \textbf{8} & 7 & 7 & 7 & 4 & 3 \\
\bottomrule
\end{tabular}
\end{table}

Table~\ref{tab:libero10} shows that \methodname{} achieves the \textbf{highest average success rate} (34.1\%) across all 10 LIBERO-10 tasks, outperforming DINOv2 by $+22\%$ relative (28.0\%), CLIP by $+56\%$ (21.9\%), SigLIP by $+26\%$ (27.1\%), and both embodied baselines by a large margin (MCR 7.8\%, R3M 4.2\%). Crucially, \methodname{} dominates on the most \emph{physics-intensive} tasks: on Task~1 (grasping a book and placing it in a compartment), \methodname{} achieves \textbf{40.0\%} vs.\ 16.0\% for DINOv2 ($+150\%$), where force modulation during grasping and placement depends on understanding object mass and contact dynamics. On Task~4 (turning on a stove and placing a pot), \methodname{} leads with 48.0\% vs.\ 32.0\% for DINOv2 ($+50\%$), and on Task~5 (\mbox{multi-object} placement), 57.3\% vs.\ 44.7\% ($+28\%$). \methodname{} is the only backbone to succeed on the hardest dual-object task (Task~0, 0.7\%) and achieves success on 8 of 10 tasks versus 7 for all other ViT baselines. Tasks~7 and~8 are universally unsolved (0\% for all backbones), representing genuine difficulty frontiers. Notably, R3M---a robotics-specific ResNet-50 pre-trained on Ego4D---achieves only 4.2\% average, confirming that egocentric video pre-training alone does not provide physics-grounded features for precise manipulation.



% ======================================================================
\section{Analysis and Discussion}
\label{sec:discussion}
% ======================================================================

\paragraph{Linear probing reveals physics-aware structure.}
Table~\ref{tab:linear_probing} shows that \methodname{} substantially outperforms all six frozen baselines on physics linear probing, achieving $R^2 = 0.550$ on mass (+46.7\% over DINOv2-B/14), $R^2 = 0.416$ on friction (+148.7\%), and $R^2 = 0.659$ on restitution (+18.5\%). Notably, \methodname{} (ViT-B/14, 86.6M params) outperforms the much larger DINOv2-L/14 (304M params) across all metrics, demonstrating that physics-aware fine-tuning is more effective than simply scaling model size. Robot-specific representations (R3M and VIP) underperform general-purpose encoders despite being pre-trained on manipulation videos, suggesting that task-agnostic visual pre-training does not capture latent physics without explicit physics-grounded supervision.

\paragraph{$k$-NN retrieval is dominated by category structure.}
Table~\ref{tab:knn} reveals that $k$-NN performance is remarkably similar across all models with strong category recognition---differences are in the third decimal place. The strict zero-shot retrieval protocol (Table~\ref{tab:zeroshot}) further confirms this: DINOv2-L/14 leads with mean $R^2{=}0.637$, while \methodname{} achieves $0.533$. This is expected: since physics properties are assigned per-category, any model that correctly identifies an object's category will retrieve neighbors with similar physics. The key insight is that \methodname{}'s advantage lies not in retrieval but in the \emph{linear decodability} of physics properties from its features (Table~\ref{tab:linear_probing}).

\paragraph{Invisible physics reveals true physics awareness.}
Table~\ref{tab:invisible_physics} provides the most compelling evidence for \methodname{}'s physics awareness. On the invisible physics test---where the same image is associated with different physics configurations---\methodname{} is the \emph{only} model with positive mass regression $R^2$ ($0.300$), while all baselines produce negative $R^2$ (worse than mean prediction). This demonstrates that \methodname{}'s contrastive training has imbued the embedding space with genuine physics-correlated structure, not merely visual-category proxies. Baseline encoders achieve near-chance accuracy (0.40--0.59) on the heavier classification task, while \methodname{} achieves 0.560 accuracy with AUC $= 0.615$, indicating above-chance sensitivity to physics-relevant features.

\paragraph{Material clustering.}
Table~\ref{tab:clustering} demonstrates that \methodname{} produces the most material-aware feature space, achieving NMI $= 0.435$ and ARI $= 0.335$---substantially above the next best baseline (CLIP-L/14: NMI $= 0.350$, ARI $= 0.314$). This confirms that the contrastive fine-tuning successfully reorganises the embedding space to group objects by material type rather than purely by semantic category.

\paragraph{Why category-grounded priors matter.}
The ablation comparing category-grounded vs.\ random physics assignment (Table~\ref{tab:ablations}) shows the largest degradation of any ablation: mass $R^2$ drops from 0.714 to 0.466 ($-34.7\%$), confirming that physically plausible training signals are essential. Without category-grounded priors, the contrastive loss cannot learn the visual-physics correlation.

\paragraph{On hard negative mining.}
An interesting finding is that removing hard negative mining actually \emph{improves} short-horizon performance (Table~\ref{tab:ablations}, 10K steps). We hypothesize that category-grounded material priors already provide strong contrastive signal: randomly sampled cross-category pairs span different materials and thus generate informative physics contrasts without explicit mining. Hard mining introduces same-category pairs with subtle within-material physics differences, which may add optimization difficulty without proportional benefit at this training scale. The full model's advantage emerges in the absolute evaluation (Table~\ref{tab:linear_probing}), where 30K steps of training with the complete pipeline achieves the best downstream performance.

\paragraph{Loss function comparison.}
The expanded ablation (Table~\ref{tab:ablations}) reveals a clear hierarchy among contrastive loss functions. Soft InfoNCE (our default) and Standard InfoNCE perform comparably, while Triplet loss degrades substantially (mass $R^2$: 0.435 vs.\ 0.601 for InfoNCE), confirming that pair-based margin losses lack the expressiveness of the full InfoNCE formulation. BYOL loss collapses entirely (mass $R^2$: 0.217, category: 72.9\%), demonstrating that the self-supervised BYOL objective---which relies on positive-only learning---cannot leverage our physics-informed contrastive pairs.

\paragraph{Fine-tuning depth.}
The Last-2 and Last-4 ablations reveal that partial fine-tuning is remarkably effective: fine-tuning only the last 4 blocks (29.9M trainable of 88.2M) achieves mass $R^2{=}0.625$, closely approaching the full fine-tuning result ($0.584$ at 30K steps, $0.714$ for no-hard-mining at 10K). This suggests that physics-relevant features can be learned primarily in the later transformer blocks, with earlier layers providing adequate general visual features from the DINOv2 pre-training.

\paragraph{Backbone fine-tuning.}
The frozen-backbone ablation shows competitive but lower physics prediction ($R^2 = 0.542$ vs.\ $0.714$ for fine-tuned), confirming that fine-tuning reshapes features beyond what a projection head alone can achieve.

\paragraph{Computational efficiency.}
Table~\ref{tab:cost} confirms that \methodname{}'s physics awareness comes at negligible computational cost. Built on DINOv2-B/14 with only a 512-dim projection head, \methodname{} adds 1.6M parameters (88.2M vs.\ 86.6M) and incurs identical latency (5.6\,ms) and throughput (873 vs.\ 874 img/s). Meanwhile, scaling to the ``large'' architecture variants (DINOv2-L/14: 304M, CLIP-L/14: 304M) triples the parameter count and halves throughput, yet these larger models are \emph{outperformed} by \methodname{} on physics probing. This highlights a key design principle: physics-relevant training signals are more valuable than additional model capacity.

\paragraph{World model dynamics.}
Table~\ref{tab:worldmodel} reveals that \methodname{} features support strong next-state dynamics prediction (mass $R^2{=}0.737$, friction $R^2{=}0.676$, restitution $R^2{=}0.807$), closely matching CLIP-L/14 despite using a 512-dim embedding vs.\ 768-dim. Interestingly, CLIP-L/14 slightly outperforms on this task, suggesting that its richer language-aligned representation provides complementary dynamics cues. However, DynaCLIP's advantage on the more challenging linear probing and invisible physics tests (Tables~\ref{tab:linear_probing},~\ref{tab:invisible_physics}) confirms that its representations encode physics structure more directly---without requiring the sequential context provided by the RSSM.

\paragraph{Physics-robust downstream policies.}
Table~\ref{tab:downstream_policy} provides compelling evidence that \methodname{}'s physics-aware features produce the most physics-robust manipulation policies. While CLIP achieves the highest absolute success rates, its performance degrades by 24\% when physics conditions change from training. In contrast, \methodname{}'s policies maintain identical performance across physics variations (generalization ratio $1.02\times$). This physics robustness aligns with \methodname{}'s design objective: by encoding physics structure directly into features, downstream policies become invariant to physics perturbations rather than relying on visual cues that correlate with specific physics at training time.

\paragraph{Transfer to established manipulation benchmarks.}
The LIBERO-10 results (Table~\ref{tab:libero10}) provide strong external validation that \methodname{}'s physics-aware features generalize beyond our custom DynaCLIP-Bench. \methodname{} achieves \textbf{34.1\%} average success rate across 10 diverse long-horizon tasks---the highest of all six backbones---outperforming DINOv2 (28.0\%, $+22\%$), CLIP-L/14 (21.9\%, $+56\%$), SigLIP (27.1\%, $+26\%$), and both robot-specific baselines by large margins (MCR 7.8\%, R3M 4.2\%). The detailed per-task breakdown reveals that \methodname{}'s advantage concentrates on \emph{physics-intensive} tasks: on ``pick up the book and place it in the compartment'' (Task~1), \methodname{} achieves \textbf{40.0\%} vs.\ 16.0\% for DINOv2 ($+150\%$), and on ``turn on the stove and put the pot on it'' (Task~4), 48.0\% vs.\ 32.0\% ($+50\%$). These tasks require precise force modulation for grasping, lifting, and placing, exactly the capabilities physics-contrastive training is designed to enhance. \methodname{} is also the only backbone to achieve any success on the hardest dual-object task (Task~0, 0.7\%) and succeeds on 8 of 10 tasks---the most of any backbone. Strikingly, R3M---a robot-specific ResNet-50 encoder---achieves success on only 3 of 10 tasks (4.2\% avg), confirming that egocentric video pre-training does not provide the physics-grounded features needed for precise manipulation.

\paragraph{Training duration.}
Extending training from 30K to 100K steps reveals an informative trade-off. Material clustering quality improves substantially (NMI: $0.435 \to 0.517$, ARI: $0.335 \to 0.480$), indicating that longer contrastive training tightens material-level cluster structure. However, linear probing $R^2$ degrades slightly (mass: $0.550 \to 0.536$, friction: $0.416 \to 0.378$, restitution: $0.659 \to 0.623$), while $k$-NN performance remains stable. We attribute this to a representation compression effect: prolonged contrastive training pulls same-material samples into tighter clusters, improving unsupervised grouping, but smooths over fine-grained \emph{within-material} physics variations that linear probes exploit. This suggests 30K steps strikes the optimal balance for downstream physics probing, while 100K steps may be preferred for clustering-oriented tasks.

\paragraph{Limitations.}
Our physics engine is analytical and deterministic, lacking collision dynamics, articulated-body effects, and deformable materials. The category-to-material mapping is coarse-grained: a \texttt{wine\_bottle} could be glass, plastic, or even metal. Category labels are limited to DomainNet's 345 classes. The $k$-NN results suggest that the learned physics structure is tightly coupled with category identity, limiting zero-shot transfer to novel categories. Scaling to open-vocabulary categories and integrating learned physics simulators are natural next steps.

\paragraph{Future Work.}
We envision three extensions: (i) replacing the analytical engine with a learned differentiable simulator~\cite{toussaint2018differentiable} for richer dynamics signals; (ii) extending to articulated and deformable objects; and (iii) integrating \methodname{} features into more sophisticated manipulation policies (\eg, Diffusion Policy~\cite{chi2023diffusion} or ACT~\cite{zhao2023act}) to further exploit physics-robust representations for contact-rich tasks.


% ======================================================================
\section{Conclusion}
\label{sec:conclusion}
% ======================================================================

We presented \methodname{}, a contrastive pre-training framework that teaches visual encoders to capture latent physical properties of objects. By leveraging category-grounded material priors and a Soft InfoNCE loss with continuous dynamics similarity targets, \methodname{} produces visual representations with physics-aware structure that is absent in general-purpose encoders. On our DynaCLIP-Bench evaluation suite of ten experiments against seven baselines (DINOv2-B/14, DINOv2-L/14, CLIP-L/14, SigLIP-B/16, R3M, VIP, MCR), \methodname{} (ViT-B/14, 88.2M params) achieves $R^2 = 0.550 / 0.416 / 0.659$ for mass/friction/restitution linear probing---outperforming DINOv2-L/14 (304M params) by 25\%/40\%/9\% and MCR (MAE) by $10.8\times$/$3.2\times$/$1.9\times$ with comparable inference latency---and leads material clustering with NMI $= 0.435$ vs.\ $0.350$ for the next best baseline. On the invisible physics test, \methodname{} is the \emph{only} model to achieve positive mass regression $R^2$ ($0.300$ vs.\ $\le -0.018$ for all baselines), demonstrating genuine physics-grounded representation rather than category proxies. On downstream manipulation in ManiSkill3, \methodname{} features produce the \emph{most physics-robust} behavior cloning policies (generalization ratio $1.02\times$ vs.\ $0.76\times$ for CLIP), and on established manipulation benchmarks, \methodname{} achieves \textbf{34.1\%} average success on LIBERO-10 ($+22\%$ over DINOv2, $+150\%$ on the most physics-intensive task). A world model trained on real ManiSkill3 trajectories achieves mass prediction $R^2{=}0.922$, confirming that mass is the most dynamically observable property. Comprehensive ablation studies over 16 variants confirm that DINOv2 initialization is essential ($-97\%$ mass $R^2$ without it), category-grounded physics priors provide $+53\%$ over random assignment, and physics-contrastive training signals outperform model scaling by $3.5\times$. We hope this work opens a new direction at the intersection of visual representation learning and physical scene understanding.


% ======================================================================
% References
% ======================================================================
\bibliographystyle{plain}
\begin{thebibliography}{99}

\bibitem{chen2020simclr}
T.~Chen, S.~Kornblith, M.~Norouzi, and G.~Hinton.
\newblock A Simple Framework for Contrastive Learning of Visual Representations.
\newblock In \emph{ICML}, 2020.

\bibitem{chen2023flatnce}
J.~Chen, M.~Hu, W.~Li, and M.~Song.
\newblock FlatNCE: Flat contrastive learning in a non-degenerate way.
\newblock \emph{arXiv preprint arXiv:2107.01152}, 2023.

\bibitem{chi2023diffusion}
C.~Chi, S.~Feng, Y.~Du, Z.~Xu, E.~Cousineau, B.~Burchfiel, and S.~Song.
\newblock Diffusion Policy: Visuomotor Policy Learning via Action Diffusion.
\newblock In \emph{RSS}, 2023.

\bibitem{guan2022neurofluid}
S.~Guan, V.~Deng, Y.~Zhao, G.~Qian, and H.~Zhao.
\newblock NeuroFluid: Fluid Dynamics Grounding with Particle-Driven Neural Radiance Fields.
\newblock In \emph{ICML}, 2022.

\bibitem{gafner2023dreamerv3}
D.~Hafner, J.~Pasukonis, J.~Ba, and T.~Lillicrap.
\newblock Mastering Diverse Domains through World Models.
\newblock \emph{arXiv preprint 

\bibitem{he2022mae}
K.~He, X.~Chen, S.~Xie, Y.~Li, P.~Doll{\'a}r, and R.~Girshick.
\newblock Masked Autoencoders Are Scalable Vision Learners.
\newblock In \emph{CVPR}, 2022.arXiv:2301.04104}, 2023.

\bibitem{hu2023maniskill3}
J.~Gu, F.~Xiang, X.~Li, Z.~Ling, X.~Liu, T.~Mu, Y.~Tang, S.~Tao, X.~Wei, Y.~Yao, X.~Yuan, P.~Xie, Z.~Huang, R.~Chen, and H.~Su.
\newblock ManiSkill3: GPU Parallelized Robotics Simulation and Benchmarking at Scale.
\newblock \emph{arXiv preprint arXiv:2410.00425}, 2024.

\bibitem{guen2020disentangling}
V.~Le~Guen and N.~Thome.
\newblock Disentangling Physical Dynamics from Unknown Factors for Unsupervised Video Prediction.
\newblock In \emph{CVPR}, 2020.

\bibitem{hafner2019dreamer}
D.~Hafner, T.~Lillicrap, I.~Fischer, R.~Villegas, D.~Ha, H.~Lee, and J.~Davidson.
\newblock Learning Latent Dynamics for Planning from Pixels.
\newblock In \emph{ICML}, 2019.

\bibitem{he2020moco}
K.~He, H.~Fan, Y.~Wu, S.~Xie, and R.~Girber.
\newblock Momentum Contrast for Unsupervised Visual Representation Learning.
\newblock In \emph{CVPR}, 2020.

\bibitem{karamcheti2023voltron}
S.~Karamcheti, S.~Nair, A.~S. Chen, T.~Kollar, C.~Finn, D.~Sadigh, and P.~Liang.
\newblock Language-Driven Representation Learning for Robotics.
\newblock In \emph{RSS}, 2023.

\bibitem{khosla2020supervised}
P.~Khosla, P.~Teterwak, C.~Wang, A.~Sarna, Y.~Tian, P.~Isola, A.~Maschinot, C.~Liu, and D.~Krishnan.
\newblock Supervised Contrastive Learning.
\newblock In \emph{NeurIPS}, 2020.

\bibitem{lerer2016learning}
A.~Lerer, S.~Gross, and R.~Fergus.
\newblock Learning Physical Intuition of Block Towers by Example.
\newblock In \emph{ICML}, 2016.

\bibitem{liu2024libero}
B.~Liu, Y.~Zhu, C.~Gao, Y.~Feng, Q.~Liu, Y.~Zhu, and P.~Stone.
\newblock LIBERO: Benchmarking Knowledge Transfer for Lifelong Robot Learning.
\newblock In \emph{NeurIPS}, 2024.

\bibitem{loshchilov2019adamw}
I.~Loshchilov and F.~Hutter.
\newblock Decoupled Weight Decay Regularization.
\newblock In \emph{ICLR}, 2019.

\bibitem{ma2022vip}
Y.~J. Ma, S.~Sodhani, D.~Jayaraman, O.~Bastani, V.~Kumar, and A.~Zhang.
\newblock VIP: Towards Universal Visual Reward and Representation via Value-Implicit Pre-Training.
\newblock In \emph{ICLR}, 2023.

\bibitem{nair2022r3m}
S.~Nair, A.~Rajeswaran, V.~Kumar, C.~Finn, and A.~Gupta.
\newblock R3M: A Universal Visual Representation for Robot Manipulation.
\newblock In \emph{CoRL}, 2022.

\bibitem{oquab2024dinov2}
M.~Oquab, T.~Darcet, T.~Moutakanni, H.~Vo, M.~Szafraniec, V.~Khalidov, P.~Fernandez, D.~Haziza, F.~Massa, A.~El-Nouby, \etal.
\newblock DINOv2: Learning Robust Visual Features without Supervision.
\newblock \emph{TMLR}, 2024.

\bibitem{peng2019domainnet}
X.~Peng, Q.~Bai, X.~Xia, Z.~Huang, K.~Saenko, and B.~Wang.
\newblock Moment Matching for Multi-Source Domain Adaptation.
\newblock In \emph{ICCV}, 2019.

\bibitem{radford2021clip}
A.~Radford, J.~W. Kim, C.~Hallacy, A.~Ramesh, G.~Goh, S.~Agarwal, G.~Sastry, A.~Askell, P.~Mishkin, J.~Clark, \etal.
\newblock Learning Transferable Visual Models From Natural Language Supervision.
\newblock In \emph{ICML}, 2021.

\bibitem{radosavovic2023vc1}
I.~Radosavovic, T.~Xiao, S.~James, P.~Abbeel, J.~Malik, and T.~Darrell.
\newblock Real-World Robot Learning with Masked Visual Pre-training.
\newblock In \emph{CoRL}, 2023.

\bibitem{raissi2019pinn}
M.~Raissi, P.~Perdikaris, and G.~E. Karniadakis.
\newblock Physics-informed neural networks: A deep learning framework for solving forward and inverse problems involving nonlinear partial differential equations.
\newblock \emph{Journal of Computational Physics}, 2019.

\bibitem{toussaint2018differentiable}
M.~Toussaint, K.~R. Allen, K.~A. Smith, and J.~B. Tenenbaum.
\newblock Differentiable physics and stable modes for tool-use and manipulation planning.
\newblock In \emph{RSS}, 2018.

\bibitem{zhai2023siglip}
X.~Zhai, B.~Mustafa, A.~Kolesnikov, and L.~Beyer.
\newblock Sigmoid Loss for Language Image Pre-Training.
\newblock In \emph{ICCV}, 2023.

\bibitem{zhu2020robosuite}
Y.~Zhu, J.~Wong, A.~Mandlekar, R.~Mart{\'\i}n-Mart{\'\i}n, A.~Joshi, S.~Nasiriany, and Y.~Zhu.
\newblock robosuite: A Modular Simulation Framework and Benchmark for Robot Learning.
\newblock \emph{arXiv preprint arXiv:2009.12293}, 2020.

\bibitem{zhu2020robosuite}
Y.~Zhu, J.~Wong, A.~Mandlekar, R.~Mart{\'i}n-Mart{\'i}n, A.~Joshi, S.~Nasiriany, and Y.~Zhu.
\newblock robosuite: A Modular Simulation Framework and Benchmark for Robot Learning.
\newblock \emph{arXiv preprint arXiv:2009.12293}, 2020.

\bibitem{zhu2020robosuite}
Y.~Zhu, J.~Wong, A.~Mandlekar, R.~Mart{\'i}n-Mart{\'i}n, A.~Joshi, S.~Nasiriany, and Y.~Zhu.
\newblock robosuite: A Modular Simulation Framework and Benchmark for Robot Learning.
\newblock \emph{arXiv preprint arXiv:2009.12293}, 2020.

\bibitem{zhu2020robosuite}
Y.~Zhu, J.~Wong, A.~Mandlekar, R.~Mart{\'\i}n-Mart{\'\i}n, A.~Joshi, S.~Nasiriany, and Y.~Zhu.
\newblock robosuite: A Modular Simulation Framework and Benchmark for Robot Learning.
\newblock \emph{arXiv preprint arXiv:2009.12293}, 2020.

\bibitem{zhao2023act}
T.~Z. Zhao, V.~Kumar, S.~Levine, and C.~Finn.
\newblock Learning Fine-Grained Bimanual Manipulation with Low-Cost Hardware.
\newblock In \emph{RSS}, 2023.

\bibitem{zheng2021ressl}
M.~Zheng, S.~You, F.~Huang, C.~Wang, Z.~Su, and C.-C.~J. Kuo.
\newblock ReSSL: Relational Self-Supervised Learning with Weak Augmentation.
\newblock In \emph{NeurIPS}, 2021.

\bibitem{zhu2020robosuite}
Y.~Zhu, J.~Wong, A.~Mandlekar, R.~Mart{\'i}n-Mart{\'i}n, A.~Joshi, S.~Nasiriany, and Y.~Zhu.
\newblock robosuite: A Modular Simulation Framework and Benchmark for Robot Learning.
\newblock \emph{arXiv preprint arXiv:2009.12293}, 2020.

\end{thebibliography}

\end{document}
